\scp{Nasals}{12-01}
PMP *n/*ŋ/*ñ > \it{n} in Banda.
This merger also occurs in \tsc{Ambon-Seram}
and could be used to link the two groups together.
However, as discussed in \srf{13-01-01},
this merger is widespread,
not exclusive (is found in other subgroups
and also in individual languages in subgroups that preserve *ŋ,
such as Ambelau), can easily be due to parallel development,
and has a number of other issues (such as traces of PMP *ñ in
some vowels) that make it a very weak basis for subgrouping.
Examples are given in \qf{12-01} and \qf{12-02}.
While Banda has a phonetic [ɲ] in a number of words,
we have found only one instance that might be a retention of
PMP *ñ.

\noindent{\wg\st{6pt}\tblr{>{\hspace{-\tabcolsep}}l<{\hspace{-\tabcolsep}}
>{\hspace{-\tabcolsep}\enspace}
lll}
\lsptoprule
%\tbx{12--1}	&	PMP	&	Banda	&	Banda gloss	\\	\midrule
%\endfirsthead
&	PMP	&	Banda	&	Banda gloss	\\	\midrule
%\endhead
\tbx{12-01}		&	*\tbr{ŋ}ajan	&	\tbr{n}alan	&	name	\\
	&	*\tbr{ŋ}ijuŋ	&	\tbr{n}ilu-, \tbr{n}ulu-n	&	nose	\\
	&	*ha\tbr{ŋ}in	&	a\tbr{n}in	&	wind	\\
	&	*na\tbr{ŋ}uy	&	na\tbr{n}o	&	swim	\\
	&	*sa\tbr{ŋ}a	&	sa\tbr{n}a-, se\tbr{n}a-no	&	branch	\\
	&	*sa\tbr{ŋ}əlaR	&	ndese\tbr{n}lar	&	fry	\\
	&	*bubu\tbr{ŋ}	&	fafu\tbr{n}	&	gable	\\
	&	*diŋdi\tbr{ŋ}	&	rindi\tbr{n}	&	cold	\\
	&	*uda\tbr{ŋ}	&	sugura\tbr{n}	&	shrimp, prawn	\\
	&	*qajə\tbr{ŋ}	&	‹\it{a}lo\tbr{n}›{\footnotemark}	&	charcoal	\\	\midrule
\tbx{12-02}	&	*\tbr{ñ}aRa	&	\tbr{n}ara-n	&	mothers sibling	\\
	&	*\tbr{ñ}awa	&	‹\tbr{nj}o\it{wa}na› &	breathe	(‹nj› = [ɲ]?) \\
	&	*pə\tbr{ñ}u	&	e\tbr{n}u	&	tortoise	\\
\lspbottomrule
\end{tabular}\mbox{}\\}

\footnotetext{Throughout this chapter examples enclosed in angle
brackets ‹ › are from \citet{vanEijbergen1864}
and are presented in the orthography used in that source.
\citet{vanEijbergen1864} follows Dutch orthographic conventions (e.g.
‹j› = [y], ‹oe› = [u])
and uses italics to indicate stress.}